\begin{problem}[33届物理竞赛决赛]{66}
    三、火星大气可视为仅由很稀薄的 $\mathrm{CO}_{2}$ 组成，此大气的摩尔质量记为 $\mu$ ，且同一高度的大气可视为处于平衡态的理想气体。火星质量为 $M_{m}$ （远大于火星大气总质量），半径为 $R_{m}$ 。假设火星大气的质量密度与距离火星表面的高度 $h$ 的关系为
    \begin{equation*}
        \rho(h) = \rho_{0} \left( 1 + \frac{h}{R_{m}} \right)^{1-n},
    \end{equation*}
    其中 $\rho_{0}$ 为常量，$n$ ($n>4$) 为常数。
    \begin{subproblem}
        求火星大气温度 $T(h)$ 随高度 $h$ 变化的表达式。
    \end{subproblem}
    \begin{subproblem}
        为了对火星表面进行探测，需将一质量为 $m_{t}$ 、质量密度较大的着陆器释放到火星表面。在某一远小于 $R_{m}$ 的高度处将该着陆器由静止开始释放，经过时间 $t_{l}$ ，着陆器落到火星表面。在着陆器下降的过程中，着陆器没有转动，火星的重力加速度和大气密度均可视为与它们在火星表面的值相等。当着陆器下降速度大小为 $v$ 时，它受到的大气阻力正比于大气密度和 $v^{2}$ 的乘积，比例系数为常量 $k$。求着陆器在着陆前的瞬间速度的大小。
    \end{subproblem}
\end{problem}